\section{Nonparametric Modelling}

\subsection{Distances}

If \(\mathcal{X} = \mathbb{R}^d\), the Minkowski distance of order \(p > 0\) between \(\vec{x}, \vec{y} \in \mathcal{X}\) is defined as:

\[\mathrm{Dis}_p(\vec{x}, \vec{y}) = \left(\sum_{j = 1}^d |x_j - y_j|^p\right)^{1 / p} = \lVert \vec{x} - \vec{y}\rVert_p\]
where, for some \(\vec{z} \in \mathcal{X}, \lVert \vec{z} \rVert_p = \left(\sum_{j = 1}^d |z_j|^p\right)^{1/p}\) is the \(p\)-norm of the vector \(\vec{z}\).

\textbf{Theorem:} The arithmetic mean minimises squared Euclidean distance. The arithmetic mean \(\mu\) of a set of data points \(D\) in a Euclidean space is the unique point that minimises the sum of squared Euclidean distances to those data points.

\subsection{Nearest Neighbour}

Stores all training examples \((x_i, f(x_i))\).

Nearest neighbour:
\begin{itemize}
    \item Given query instance \(x_q\), first locate nearest training example \(x_n\), then estimate \(\hat{f}(x_q) \leftarrow f(x_n)\).
\end{itemize}

\(k\)-Nearest neighbour:
\begin{itemize}
    \item Given \(x_q\), take vote among its \(k\) nearest neighbours (if discrete-valued target function)
    \item take mean of \(f\) values of \(k\) nearest neighbours (if real-valued)
    \[\hat{f}(x_q) \leftarrow \frac{\sum_{i = 1}^k f(x_i)}{k}\]
\end{itemize}

\subsubsection{\(k\)-Nearest Neighbour (\(k\)NN) Algorithm}

Training algorithm:
\begin{itemize}
    \item For each training example \((x_i, f(x_i))\), add the example to the list \textit{training\_examples}.
\end{itemize}

Classification algorithm:
\begin{itemize}
    \item Given a query instance \(x_q\) to be classified,
    \begin{itemize}
        \item Let \(x_1 \dots x_k\) be the \(k\) instances from \textit{training\_examples} that are nearest to \(x_q\) by the distance function
        \item Return
        \[\hat{f}(x_q) \leftarrow \argmax_{y \in \mathcal{Y}} \sum_{i = 1}^k I[y, f(x_i)]\]
        where \(I[a, b] = 1\) if \(a = b\) and 0 otherwise.
    \end{itemize}
\end{itemize}

Use Nearest Neighbour when
\begin{itemize}
    \item Instances map to points in \(\mathbb{R}^d\)
    \item Low-dimensional data, say less than 20 features per instance (Curse of Dimensionality)
    \item Lots of training data
    \item Target function (classification or regression) is ``local''
    \item No requirement for ``explanatory'' model to be learned
    \item Small \(k\) has low bias, high variance. Increasing \(k\) increases bias and decreases variance.
\end{itemize}

\subsection{Tree Learning}

\subsubsection{Top-Down Induction of Decision Trees (TDIDT)}

\begin{enumerate}
    \item Select \(A\) as the ``best'' decision attribute for the next \textit{node}
    \item Assign \(A\) as decision attribute for \textit{node}
    \item For each value of \(A\), create new descendant of \textit{node}
    \item Sort training examples to leaf nodes
    \item If training examples perfectly classified, then STOP, else iterate over new leaf nodes
\end{enumerate}

\subsubsection{Information Gain}

Entropy measures the ``impurity'' of \(S\). \(Entropy(S) = \) expected number of bits needed to encode class (\(\oplus\) or \(\ominus\)) of randomly drawn member of \(S\) (under the optimal, shortest-length code).

\[Entropy(S) = -p_\oplus \log_2 p_\oplus - p_\ominus \log_2 p_\ominus\]

\(Gain(S, A) = \) expected reduction in entropy due to sorting on \(A\).

\[Gain(S, A) \equiv Entropy(S) - \sum_{v \in Values(A)} \frac{|S_v|}{|S|}Entropy(S_v)\]

\subsubsection{Overfitting}

To avoid overfitting in Tree learning there are two main approaches
\begin{itemize}
    \item \textbf{pre-pruning}: stop growing when further data splits are not useful
    \item \textbf{post-pruning} grow full tree, then remove sub-trees which may be overfitting
\end{itemize}

Post-pruning more common:
\begin{itemize}
    \item can prune using a pessimistic estimate of error
    \begin{itemize}
        \item measure training set error
        \item adjust by a factor dependent on ``confidence'' hyperparameter
    \end{itemize}
    \item or can use cross-validation to estimate error during pruning
\end{itemize}

\subsubsection{Regression Trees}

Difference to decision trees:
\begin{itemize}
    \item Splitting criterion: minimising intra-subset variation
    \item Pruning criterion: based on numeric error measure
    \item Leaf node predicts average class values of training instances reaching that node
\end{itemize}

Can approximate piecewise constant functions, easy to interpret, more sophisticated version: model trees.

% tree learning as variance reduction slide 63

\subsubsection{Model Trees}
\begin{itemize}
    \item Like regression trees but with linear regression functions at each node
    \item Linear regression applied to instances that reach a node after full tree has been built
    \item Only a subset of the attributes is used for linear regression. Attributes occurring in subtree (+ maybe attributes occurring in path to the root)
    \item Fast: overhead for linear regression not large because usually only a small subset of attributes is used in tree
\end{itemize}